% ============================================================================
% CHAPITRE 5 — À INSÉRER APRÈS « 5.10 Filtrage par domaine et spécialité »
% ============================================================================
% Ce fichier remplace les sections 5.11 (Génération contrôlée) et 5.12
% (Optimisation et intégration) de la version actuelle du mémoire.
%
% Nouvelle numérotation obtenue :
%   5.11  Évaluation du pipeline et limites identifiées      (nouveau)
%   5.12  Traitement déterministe des questions temporelles  (nouveau)
%   5.13  Génération contrôlée                               (ex-5.11, enrichie)
%   5.14  Optimisation et intégration                        (ex-5.12, enrichie)
%
% Le tableau 5.6 « Comparaison des fournisseurs de modèle de langage » de la
% version actuelle est conservé tel quel en 5.13.1 : il n'est pas reproduit ici.
% ============================================================================


% ============================================================================
\section{Évaluation du pipeline et limites identifiées}
\label{sec:eval-rag}
% ============================================================================

Après la mise en œuvre du pipeline RAG initial, des tests ont été réalisés sur
différentes catégories de questions. Ces tests ont permis d'identifier certaines
limites qui ne pouvaient pas être entièrement résolues par la seule recherche
documentaire. Deux cas particuliers ont notamment été observés : les questions
nécessitant une interprétation temporelle ou un calcul lié à une date, et la
nécessité d'encadrer davantage la génération de la réponse finale.

Ces deux constats ont donné lieu à deux mécanismes supplémentaires, décrits dans
les sections~\ref{sec:calendrier} et~\ref{sec:generation}, qui viennent compléter
le pipeline sans en modifier le principe : la recherche documentaire reste
inchangée, mais son résultat n'est plus transmis seul au modèle de langage.

\subsection{Protocole de test}

L'évaluation a été organisée en trois niveaux, volontairement séparés : une
défaillance de la génération et une défaillance de la recherche n'ont ni la même
cause ni le même correctif, et les confondre dans un test unique de bout en bout
rend le diagnostic impossible.

\begin{table}[h]
\centering
\caption{Niveaux de test du pipeline RAG}
\label{tab:niveaux-test}
\begin{tabular}{|p{4.1cm}|p{5.4cm}|p{4.1cm}|}
\hline
\textbf{Niveau} & \textbf{Ce qu'il vérifie} & \textbf{Dépendances} \\
\hline
Recherche seule \newline (\texttt{test\_retrieval.py},
point d'accès \texttt{/api/chat/debug})
& Les articles attendus remontent-ils dans les extraits retenus ?
Inspection qualitative des 8 questions ciblant chacune un mécanisme différent
(dense, BM25, accents, réponse composite).
& Index local uniquement, aucun appel au modèle de langage \\
\hline
Mesure quantitative \newline (\texttt{evaluer.py})
& Métriques chiffrées et reproductibles sur le jeu de questions annotées,
comparables entre deux versions du système.
& Index local ; appel au modèle uniquement pour la partie génération \\
\hline
Calcul des dates \newline (\texttt{test\_calendrier.py})
& Le verdict calculé est-il conforme à l'arrêté pour 9 cas couvrant chacun
une règle distincte ?
& Aucune : ni réseau, ni modèle d'embedding \\
\hline
\end{tabular}
\end{table}

\paragraph{Jeu de questions annotées.}
La mesure quantitative suppose une référence. Un fichier
\texttt{data/eval/gold\_standard.json} a donc été constitué : pour chaque
question, il associe la liste des identifiants de chunks qui \emph{devraient}
être retrouvés — vérifiés un à un dans l'index — et une réponse de référence
rédigée à la main à partir du texte officiel. Il compte actuellement
7~questions du domaine \texttt{chasse}, chacune documentant le mécanisme
qu'elle cible : recherche dense lorsque la question ne partage aucun mot avec
l'article visé, recherche lexicale lorsqu'elle contient un terme rare
(« adjudicataires »), normalisation des accents (« Aïn Draham »), ou réponse
composite exigeant deux articles distincts.

\paragraph{Métriques retenues.}
Les métriques de recherche sont calculées sur les $k = 8$ extraits réellement
transmis au modèle, et non sur une valeur de $k$ choisie pour l'occasion : la
mesure porte ainsi sur ce que le système fait, et non sur une configuration
théorique.

\begin{table}[h]
\centering
\caption{Métriques d'évaluation du pipeline}
\label{tab:metriques}
\begin{tabular}{|p{3.3cm}|p{5.6cm}|p{4.7cm}|}
\hline
\textbf{Métrique} & \textbf{Ce qu'elle mesure} & \textbf{Principe de calcul} \\
\hline
Precision@8 & Parmi les 8 extraits retournés, la proportion réellement
pertinente & extraits pertinents retrouvés / 8 \\
\hline
Recall@8 & Parmi les extraits pertinents existants, la proportion présente
dans les 8 retenus & extraits pertinents retrouvés / total attendu \\
\hline
MRR & Le rang auquel apparaît le premier extrait pertinent & moyenne de
$1/\text{rang du premier extrait pertinent}$ \\
\hline
nDCG@8 & Comme le rappel, mais pénalise un extrait pertinent trouvé au rang 8
plutôt qu'au rang 1 & somme pondérée par la position, normalisée \\
\hline
Fidélité & La réponse s'appuie-t-elle uniquement sur les extraits fournis,
sans information ajoutée ? & un second appel au modèle juge la réponse au
regard du contexte transmis \\
\hline
Pertinence & La réponse traite-t-elle réellement la question posée, sans
digression ? & même principe, sur le couple question / réponse \\
\hline
Latence & Temps de réponse effectif, moyenne et 95\textsuperscript{e}
centile & mesuré par requête et agrégé sur le jeu de test \\
\hline
\end{tabular}
\end{table}

Les deux dernières métriques de fidélité et de pertinence reprennent le principe
du framework \textsc{ragas}, réimplémenté directement dans le script
d'évaluation : un second appel au même modèle juge la sortie du premier et
renvoie un score normalisé accompagné d'une justification. Ce choix évite une
dépendance supplémentaire et garde le critère d'évaluation lisible dans le code.
Lorsque le jugement ne peut pas être exploité, le score est écarté du calcul
plutôt que compté comme nul, afin de ne pas dégrader la moyenne par un incident
technique.

\begin{table}[h]
\centering
\caption{Résultats de l'évaluation (jeu de 7 questions annotées, $k = 8$)}
\label{tab:resultats-eval}
\begin{tabular}{|p{3.3cm}|p{2.6cm}|p{7.7cm}|}
\hline
\textbf{Métrique} & \textbf{Valeur} & \textbf{Lecture} \\
\hline
% À COMPLÉTER après exécution de : python -m scripts.evaluer
% Les valeurs sont écrites dans data/eval/resultats.json
Precision@8 & \dots & \dots \\
\hline
Recall@8 & \dots & \dots \\
\hline
MRR & \dots & \dots \\
\hline
nDCG@8 & \dots & \dots \\
\hline
Fidélité & \dots & \dots \\
\hline
Pertinence & \dots & \dots \\
\hline
Latence moyenne & \dots & \dots \\
\hline
\end{tabular}
\end{table}

\subsection{Limites relevées}

Le tableau~\ref{tab:lacunes} récapitule les limites constatées. Certaines ont été
observées directement dans les réponses produites lors des tests ; d'autres
relèvent d'un risque identifié en relisant ces réponses, avant qu'il ne se
manifeste sur une question réelle.

\begin{table}[h]
\centering
\caption{Limites identifiées après tests et correctifs apportés}
\label{tab:lacunes}
\begin{tabular}{|p{4.5cm}|p{4.2cm}|p{4.6cm}|}
\hline
\textbf{Constat} & \textbf{Origine} & \textbf{Correctif} \\
\hline
Sur une question de date précise, la réponse est plausible mais le
raisonnement est faux : le modèle a écrit « demain est un mercredi » pour un
mardi, en inventant au passage une contrainte d'horaire matinal
& Répondre suppose de croiser sept conditions ; le modèle en produit une
approximation vraisemblable au lieu de refuser
& Calcul du verdict en Python, transmis au modèle
(section~\ref{sec:calendrier}) \\
\hline
Une date révolue peut être présentée comme « autorisée », l'utilisateur
comprenant qu'il peut s'y rendre
& Un unique état autorisé / refusé ne distingue pas une période ouverte d'une
période écoulée
& Troisième état « date passée », indépendant de l'autorisation
(section~\ref{sec:verdict}) \\
\hline
Le verdict fourni est reformulé avec des circonstances que l'utilisateur n'a
pas mentionnées (heure, lieu, matériel)
& Le verdict était transmis comme un extrait parmi d'autres, sans statut
particulier
& Ligne « réponse à donner » placée en tête du verdict et règle interdisant
tout recalcul (section~\ref{sec:prompt}) \\
\hline
Le vocabulaire interne du système apparaît dans la réponse
& Aucune consigne ne précisait que ces blocs sont invisibles pour
l'utilisateur
& Règle de restitution : reformuler les réserves en langage courant
(section~\ref{sec:prompt}) \\
\hline
Une période d'arrêté déjà écoulée est présentée comme en vigueur
& La date du jour n'était pas transmise au modèle
& Date calculée dynamiquement à chaque appel et règle de validité temporelle
(section~\ref{sec:prompt}) \\
\hline
Réponses longues et hors sujet sur une salutation ou une question étrangère
au domaine
& Aucun cas prévu pour les messages ne comportant pas de question
réglementaire
& Deux règles de brièveté, une phrase par réponse
(section~\ref{sec:prompt}) \\
\hline
« quel est le prix du permis » ne remonte aucun montant chiffré ; « que
risque un chasseur » aucune sanction
& Écart de vocabulaire entre l'usager et le texte officiel, qui dit
« redevance domaniale » et « est puni de »
& Expansion de requête côté recherche lexicale (section~5.9.3) \\
\hline
Les réponses composites sont tronquées : le montant du permis relève de deux
articles, les sanctions de trois
& Nombre d'extraits transmis trop faible
& Passage à 8 extraits retenus, pour un surcoût de contexte négligeable \\
\hline
Une question de suivi telle que « et pour le lièvre ? » ne retrouve rien
& Prise isolément, elle ne contient aucun mot-clé exploitable
& Concaténation du tour précédent à la requête de recherche uniquement \\
\hline
\end{tabular}
\end{table}

Les trois derniers constats concernent des réglages du pipeline déjà décrit :
ils ont été corrigés par un ajustement de la recherche, sans nouveau composant.
Les six premiers, en revanche, portent tous sur ce qui se passe \emph{après} la
recherche documentaire — le traitement des dates et le contrôle de la
génération. Ils ont conduit aux deux ajouts présentés ci-après.


% ============================================================================
\section{Traitement déterministe des questions temporelles}
\label{sec:calendrier}
% ============================================================================

\subsection{Pourquoi la recherche documentaire ne suffit pas}

La question « puis-je chasser le sanglier le 21 janvier ? » n'admet pas de
réponse directement récupérable dans un article : aucun texte ne l'énonce. Y
répondre suppose de croiser sept conditions successives :

\begin{enumerate}
  \item l'espèce citée, qui détermine la période applicable ;
  \item la période d'ouverture correspondante (article premier de l'arrêté) ;
  \item le jour de la semaine, tous les jours n'étant pas ouverts (article 4) ;
  \item les jours fériés officiels, qui ouvrent des jours normalement fermés ;
  \item les exceptions géographiques, dont l'interdiction du 15 octobre 2025
        dans les gouvernorats d'Ariana et de Bizerte ;
  \item le statut du demandeur, la chasse touristique étant restreinte à
        certaines espèces (article 18) ;
  \item la date du jour, qui détermine si la saison est encore en cours.
\end{enumerate}

Les tests ont montré qu'un modèle de langage n'exécute pas cette chaîne de façon
fiable. Il en produit une approximation vraisemblable, ce qui est plus dommageable
qu'un refus : l'utilisateur ne dispose d'aucun moyen de distinguer une réponse
juste d'une réponse fausse. Le cas relevé est instructif : la conclusion était
correcte, mais obtenue par un raisonnement faux — un mardi qualifié de mercredi,
et une contrainte d'horaire inventée. Une conclusion juste par accident n'est pas
une réponse fiable, puisque rien ne garantit que le même raisonnement
produira une conclusion juste la fois suivante.

Ce calcul a donc été retiré au modèle et confié à un module Python dédié,
\texttt{rag/calendrier.py}. Le raisonnement y devient déterministe, testable et
auditable : les mêmes entrées produisent toujours le même verdict, et chaque
règle appliquée renvoie à l'article dont elle est issue.

\subsection{Données de référence}

Le module encode les neuf groupes d'espèces de l'arrêté du 27 août 2025, avec
pour chacun sa période, ses jours autorisés et ses restrictions propres. Les
sept jours fériés civils à date fixe y figurent également.

\begin{table}[h]
\centering
\caption{Périodes et jours de chasse encodés (arrêté du 27 août 2025)}
\label{tab:especes}
\begin{tabular}{|p{3.5cm}|p{3.4cm}|p{4.3cm}|p{1.6cm}|}
\hline
\textbf{Groupe d'espèces} & \textbf{Période} & \textbf{Jours autorisés} &
\textbf{Article} \\
\hline
Petit gibier sédentaire (lièvre, perdrix, caille sédentaire, pigeon biset,
ganga unibande) & 5 oct. 2025 au 7 déc. 2025 & Dimanches et jours fériés
officiels & 1 et 4 \\
\hline
Gros gibier (sanglier, loup doré africain) & 5 oct. 2025 au
1\textsuperscript{er} févr. 2026, reporté au 26 avr. 2026 dans huit
gouvernorats du sud & Tous les jours de la semaine & 1 et 4 \\
\hline
Pigeon ramier & 9 nov. 2025 au 15 mars 2026 & Du mercredi au dimanche et
jours fériés & 1 et 4 \\
\hline
Gibier d'eau & 9 nov. 2025 au 15 mars 2026 & Du mercredi au dimanche et
jours fériés & 1 et 4 \\
\hline
Grives et étourneaux & 9 nov. 2025 au 15 mars 2026 ; du 7 déc. 2025 au
15 mars 2026 pour les touristes & Du mercredi au dimanche et jours fériés &
1, 4 et 18 \\
\hline
Bécasse des bois & 9 nov. 2025 au 15 mars 2026 & Du mercredi au dimanche et
jours fériés, dans neuf gouvernorats, sans battue ni poste & 1 et 4 \\
\hline
Caille à l'épervier & 5 avr. 2026 au 28 juin 2026 & Selon l'article 4,
gouvernorat de Nabeul uniquement & 1 \\
\hline
Tourterelle de passage et sédentaire & 26 juil. 2026 au 20 sept. 2026 & Du
mercredi au samedi à partir de 15 h, toute la journée les dimanches et jours
fériés & 1 et 4 \\
\hline
Gangas (cata, tacheté, couronné) & 26 juil. 2026 au 20 sept. 2026 & Du
mercredi au dimanche et jours fériés & 1 et 4 \\
\hline
\end{tabular}
\end{table}

L'espèce est reconnue dans la question par correspondance sur une liste de
dénominations courantes, y compris les appellations locales, en testant les
expressions les plus longues d'abord : « caille à l'épervier » doit primer sur
« caille », et « pigeon ramier » sur « pigeon biset ». La date est reconnue sous
plusieurs formes — « aujourd'hui », « demain », « le 21 janvier », « 21/01/2026 ».
Lorsque l'année n'est pas précisée, celle qui place la date dans la saison en
cours ou à venir est retenue, plutôt qu'une année arbitraire.

\subsection{Le verdict et ses trois états}
\label{sec:verdict}

Le résultat du calcul n'est pas un booléen mais une structure comportant
l'autorisation, l'indication que la date est révolue, la liste des constats
ayant conduit au verdict et la liste des points que le calcul ne couvre pas.

La distinction entre « autorisé » et « passé » est un enseignement direct des
tests. Ces deux propriétés sont indépendantes : le 21 janvier 2026 tombait dans
la période d'ouverture du gros gibier, mais annoncer cette date « autorisée »
aujourd'hui laisserait entendre que l'utilisateur peut s'y rendre. Le verdict
énonce donc trois états — autorisé, non autorisé, date passée — et le troisième
prime sur les deux premiers dans la formulation transmise au modèle.

Le module produit également une phrase de conclusion en français courant, que le
modèle n'a plus qu'à reformuler. Cette phrase est placée \emph{avant} le détail
du raisonnement dans le bloc transmis, et non après : reléguée en fin de bloc,
elle se trouvait noyée dans les constats et le modèle reconstruisait son propre
raisonnement à partir de ceux-ci, avec le résultat décrit plus haut.

\subsection{Ce que le calcul ne couvre pas}

Deux points ont été délibérément laissés hors du calcul, parce qu'ils ne peuvent
pas être traités de façon fiable :

\begin{itemize}
  \item les réserves de chasse énumérées à l'article 12, qui recensent plusieurs
        centaines de toponymes dont l'orthographe varie trop pour permettre une
        correspondance sûre ;
  \item les fêtes religieuses, qui suivent le calendrier lunaire et ne se
        déduisent pas d'une date civile — or elles ouvrent des jours
        normalement fermés.
\end{itemize}

Dans les deux cas, le verdict le signale explicitement au lieu de deviner. Ces
réserves sont transmises au modèle, qui doit les restituer à l'utilisateur : une
réponse indiquant qu'il reste à vérifier si le jour est férié, ou si le lieu
figure parmi les réserves, est préférable à une réponse affirmative qui
ignorerait ces conditions. La limite du calcul devient ainsi une information
utile plutôt qu'un silence.

\subsection{Restitution à l'interface}

Le verdict est exposé sous deux formes. La première est le bloc de texte injecté
dans le prompt, destiné au modèle. La seconde est une structure sérialisable
destinée à l'application : un état, une phrase courte pour le bandeau
(« Non, pas demain »), l'espèce reconnue et, lorsque la réponse est négative, la
date de la prochaine ouverture. Cette dernière est obtenue en parcourant les
jours suivants jusqu'à en trouver un qui satisfasse à la fois la période et les
jours autorisés. L'interface peut ainsi afficher une information exploitable là
où le système ne renvoyait qu'un refus.

Le calcul des dates a par ailleurs une conséquence sur le cache de réponses :
une question portant sur une date n'y est jamais enregistrée, « demain » ne
désignant pas le même jour d'un jour à l'autre.

\subsection{Validation}

Le module est validé par un jeu de cas exécutable sans réseau ni modèle
d'embedding, ce qui le rend rejouable à chaque modification. Chaque cas vérifie
le couple (date passée, autorisation) plutôt que la seule autorisation, afin que
la distinction introduite plus haut soit elle aussi testée. Les cas sont évalués
à une date de référence fixe, faute de quoi leur résultat changerait avec le jour
d'exécution.

\begin{table}[h]
\centering
\caption{Cas de validation du calcul des dates (date de référence :
10 août 2026)}
\label{tab:cas-calendrier}
\begin{tabular}{|p{4.9cm}|p{2.4cm}|p{6.1cm}|}
\hline
\textbf{Cas} & \textbf{Attendu} & \textbf{Règle vérifiée} \\
\hline
Tourterelle, 12 août 2026 & Autorisé & Période en cours et jour de la semaine
autorisé \\
\hline
Tourterelle, 10 août 2026 & Refusé & Un lundi ne figure pas parmi les jours
ouverts (article 4) \\
\hline
Sanglier, aujourd'hui & Refusé & La saison 2025/2026 est fermée depuis
février \\
\hline
Sanglier, 21 janvier 2026 & Passée, alors ouverte & Une date révolue ne
s'annonce pas autorisée \\
\hline
Sanglier, 20 mars 2026 à Tozeur & Passée, alors ouverte & Fermeture reportée
dans les huit gouvernorats du sud \\
\hline
Sanglier, 20 mars 2026 à Bizerte & Passée et fermée & Hors de ces huit
gouvernorats, la période était déjà close \\
\hline
Touriste, tourterelle, 12 août 2026 & Refusé & Chasse touristique limitée au
gros gibier et aux grives (article 18) \\
\hline
Lièvre, 21 janvier 2026 & Passée et fermée & Le petit gibier ferme dès le
7 décembre \\
\hline
Ganga, 16 août 2026 & Autorisé & Un dimanche est ouvert et la période est en
cours \\
\hline
\end{tabular}
\end{table}


% ============================================================================
\section{Génération contrôlée}
\label{sec:generation}
% ============================================================================

Le contexte assemblé — extraits retenus par la recherche et verdict éventuel
calculé par le module précédent — est transmis au modèle de langage chargé de
formuler la réponse finale. Les tests ont montré que fournir un bon contexte ne
suffit pas : encore faut-il contraindre ce que le modèle en fait. Cette
génération est donc \emph{contrôlée} : le modèle ne dispose d'aucune liberté
d'invention, sa tâche se limite à interpréter et reformuler ce qui lui est
fourni, jamais à rechercher ou à calculer par lui-même.

\subsection{Fournisseur du modèle de langage}

% Conserver ici le tableau 5.6 de la version actuelle du mémoire
% (« Comparaison des fournisseurs de modèle de langage »), inchangé.

\subsection{Assemblage du contexte}

L'ordre dans lequel les éléments sont présentés au modèle n'est pas neutre. Le
verdict est placé avant les extraits, car il fait autorité sur la question de
date, les extraits ne servant qu'à l'étayer. Chaque extrait porte ensuite son
étiquetage complet : source, article, rang normatif et validité. C'est cet
étiquetage qui rend applicables les règles de citation, de validité temporelle et
de hiérarchie des normes énoncées ci-après — sans lui, le modèle ne pourrait ni
citer un article, ni détecter une période close, ni arbitrer entre une loi et un
arrêté.

Les tours de conversation précédents sont transmis séparément, sous forme
d'échanges classiques, et leur nombre est borné : au-delà de quelques échanges,
l'ancien contexte devient du bruit. Les extraits, eux, restent attachés au
dernier message : ils sont propres à la question courante, et il serait faux de
les rattacher aux tours antérieurs.

\subsection{Prompt système}
\label{sec:prompt}

Le comportement du modèle est encadré par onze règles numérotées. Plusieurs
d'entre elles ne relèvent pas d'une précaution théorique mais répondent
directement à un comportement observé lors des tests.

\begin{table}[h]
\centering
\caption{Règles du prompt système et lacune visée}
\label{tab:regles-prompt}
\begin{tabular}{|p{3.4cm}|p{10.3cm}|}
\hline
\textbf{Règle} & \textbf{Ce qu'elle impose, et pourquoi} \\
\hline
1. Fidélité & Répondre uniquement à partir des extraits fournis et le dire
lorsque l'information n'y figure pas, plutôt que de compléter par des
connaissances générales : c'est la règle dont la métrique de fidélité mesure le
respect \\
\hline
2. Citation & Indiquer le texte et le numéro d'article pour toute règle
juridique, de sorte que l'utilisateur puisse vérifier ; inutile pour les
questions portant sur l'application \\
\hline
3. Validité temporelle & Comparer la validité de l'extrait à la date du jour et
prévenir si la période est écoulée : la date est recalculée à chaque appel,
jamais écrite en dur, sans quoi le système annoncerait des périodes closes comme
en vigueur \\
\hline
4. Hiérarchie des normes & En cas d'apparente contradiction, la loi prime sur
l'arrêté ; mais un arrêté peut légalement prévoir une dérogation, et il faut
alors le signaler au lieu de conclure à une contradiction \\
\hline
5. Avertissements & Traiter en priorité les extraits signalant un renvoi
périmé, produits comme tels au moment du découpage du corpus \\
\hline
6. Dates & Ne jamais calculer soi-même si une date est autorisée. Lorsqu'un
verdict est fourni, sa ligne « réponse à donner » est la réponse : elle se
reformule sans changer de sens, et ni le jour de la semaine, ni la période, ni
l'espèce ne se recalculent ; aucune circonstance que l'utilisateur n'a pas
mentionnée ne doit être ajoutée. À défaut de verdict, exposer la règle et
inviter à vérifier, sans conclure à la place de l'utilisateur \\
\hline
7. Aucun jargon interne & L'utilisateur ne voit ni les extraits, ni les blocs
transmis au modèle : ces termes ne doivent jamais apparaître dans la réponse.
Une réserve se restitue en langage courant — « pensez à vérifier si ce jour est
férié » — et non par une mention du mécanisme interne \\
\hline
8. Messages sans question & Répondre par une seule phrase courte à une
salutation ou un remerciement, sans se présenter ni énumérer ses compétences \\
\hline
9. Sécurité & En cas de doute sur une question réglementaire, inviter à
contacter l'arrondissement régional des forêts \\
\hline
10. Forme & Répondre brièvement et directement, sans titres ni listes sauf
nécessité, et dans la langue de la question, y compris en arabe ou en dialecte
tunisien \\
\hline
11. Hors sujet & Répondre en une phrase et rappeler son domaine, sans décrire
le contenu des extraits ni demander de reformuler la question \\
\hline
\end{tabular}
\end{table}

La température du modèle est par ailleurs fixée à une valeur basse (0,2), qui
privilégie la fidélité au texte source plutôt que la créativité rédactionnelle,
et la longueur de la réponse est bornée. Ces deux réglages vont dans le même sens
que les règles ci-dessus : sur un corpus réglementaire, la variabilité de
formulation n'apporte rien et le risque d'écart au texte augmente avec elle.

\subsection{Sources restituées avec la réponse}

Le contrôle de la génération ne repose pas uniquement sur le prompt. La réponse
renvoyée à l'application est accompagnée des sources effectivement transmises au
modèle — texte, article, gouvernorat et validité — que l'interface affiche sous
la réponse. L'utilisateur peut ainsi remonter au texte officiel, ce qui rend une
éventuelle imprécision détectable au lieu de rester invisible. Seuls les extraits
de nature normative y figurent, et au plus quatre : les contenus d'aide sur
l'application n'ont pas de source réglementaire à citer, et une liste plus longue
que la réponse elle-même nuirait à sa lisibilité.


% ============================================================================
\section{Optimisation et intégration}
\label{sec:optimisation}
% ============================================================================

Deux derniers éléments complètent le système sans modifier son fonctionnement :
un cache de réponses, qui réduit la dépendance au réseau, et une API qui expose
l'ensemble du pipeline à l'application Ghabetna.

Le modèle d'embedding et les index sont chargés une seule fois, au démarrage du
service, plutôt qu'à chaque question : leur chargement prend plusieurs secondes,
ce qui serait incompatible avec un temps de réponse raisonnable s'il devait être
répété à chaque requête. Le service reste par ailleurs sans état : l'historique
de conversation est conservé par l'application et renvoyé à chaque appel, ce qui
permet de répliquer le service sans coordination et lui évite toute base de
données propre.

Chaque question reçue par l'API traverse alors l'ensemble du pipeline déjà
décrit — recherche, filtrage par domaine, calcul éventuel du verdict, génération
— avant que la réponse ne soit renvoyée avec ses sources et le verdict éventuel.
Un second point d'accès expose uniquement le résultat de la recherche, sans
appeler le modèle de langage : il a permis de vérifier la pertinence des extraits
retrouvés pour les nouveaux domaines camper et apiculture sans consommer de quota
d'appel, et sert de premier niveau de diagnostic lorsqu'une réponse paraît
incorrecte — il indique si le problème vient de la recherche ou de la génération.

Un cache de réponses, sous la forme d'un simple dictionnaire persistant sur
disque dont la clé combine la question normalisée et la spécialité, sert de filet
de sécurité réseau : il n'est ni un second modèle, ni une forme d'apprentissage.
Il n'est volontairement pas appliqué lorsque la question s'inscrit dans un
historique de conversation, la réponse dépendant alors du contexte précédent, ni
lorsqu'elle porte sur une date, pour la raison exposée en
section~\ref{sec:calendrier}. La spécialité entre dans la clé parce qu'un
chasseur et un apiculteur qui posent la même question interrogent des corpus
différents et doivent donc recevoir des réponses différentes. Posé une fois pour
chaque question de démonstration avant une présentation, ce cache permet au
système de continuer à répondre même en l'absence de connexion réseau.
